% Proposed replacement for the introduction in main_v2.tex.
% This file is a review draft; main_v2.tex has not been changed.
% All cited references are in bib.bib; use \bibliography{bib}.
%
% Original introduction, preserved verbatim below (each line is commented out):
% \section{Introduction to the architecture}
% Turbulent flows are known as a notoriously difficult framework for numerical simulations. One of the common method in such simulations is the \textit{direct numerical simulations} (DNS) that discretize the Navier-Stokes equation to a space and time grid and then simulate directly. Such discretizations, however require a very fine grid resolution scaling as $N_{\textrm{grid}} \sim \left(Re\right)^{\alpha}$ with $\alpha \sim 2.05 $ \cite{yang2021grid}. Insufficient small-scale resolution may leave mean dissipation and enstrophy approximately converged while substantially underestimating high-order moments of velocity gradients, dissipation, and enstrophy, as well as the tails of velocity-gradient distributions and statistics involving second velocity derivatives \cite{DonzisResolutionEffects}.
%  While this becomes expensive, a variety of models that regularize the equations on a larger scale, introducing proper dissipative terms at the cost of short time accuracy \cite{Wilson1996Review}. Another approach consists in keeping coarse-grained scales, without non-controlled dissipative approximation, but at a cost of keeping the non-Markovian memory terms \cite{YenTingLinMZ2021,de2026data}.
%
% In this article, we investigate learned memory corrections for weather prediction by augmenting a frozen GraphCast model with a trainable branch that combines spatial message passing and temporal memory.

\section{Introduction}

Many numerical and learned forecasters advance a resolved weather state from its current value or a short window of recent values. Examples include the numerical Integrated Forecasting System (IFS) configuration used to produce ERA5~\cite{ecmwf2016ifsdynamics,hersbach2020era5}, and learned Pangu-Weather~\cite{bi2023pangu}, GraphCast~\cite{lam2023graphcast}, FuXi~\cite{chen2023fuxi}, ECMWF's ERA5-native-grid AIFS~\cite{lang2024aifs}, WeatherNext 2~\cite{alet2025weathernext2}, and WeatherNext 3~\cite{rasp2026weathernext3}. This is a practical representation for autoregressive prediction, but the resolved variables need not contain every influence on their subsequent evolution. In global weather prediction, finite horizontal resolution, limited observed variables, and approximate representations of smaller-scale processes all constrain the state available to a forecast model. Whether information about omitted processes can be recovered from the history of that resolved state is therefore a question about the chosen model representation, rather than a claim that the full atmospheric dynamics lack a Markovian state.

The difficulty of resolving all relevant scales is familiar from turbulent-flow simulation. Direct numerical simulation (DNS) discretizes the governing equations finely enough to represent dynamically active scales, and its grid requirements grow rapidly with Reynolds number. For example, a grid-point estimate for a spatially developing turbulent boundary layer gives $N_{\mathrm{grid}}\sim Re_{L_x}^{2.05}$ under locally adapted DNS spacing~\cite{yang2021grid}. This is a cost estimate for that flow, not a forecast-error law for the atmosphere. Resolution also affects what a simulation can represent: insufficient small-scale resolution can leave mean dissipation and enstrophy nearly converged while strongly underestimating higher-order gradient statistics and distribution tails~\cite{DonzisResolutionEffects}. Coarser practical models must represent the effects of unresolved motions through an effective closure.

Such a closure need not depend only on the present resolved state. Projecting full dynamics onto a smaller set of variables can yield an instantaneous contribution, a contribution depending on the history of those variables, and a term associated with unobserved initial conditions. The Mori--Zwanzig formalism makes this decomposition explicit; data-driven studies have learned its memory contribution in reduced dynamical systems~\cite{YenTingLinMZ2021,de2026data}. This motivates temporal memory as a possible way to represent effects left out by a finite-window predictor. It does not imply that a learned latent state can be identified with a unique physical memory kernel.

For the Kuramoto--Sivashinsky equation, the advantage of memory over a memoryless operator changes with grid resolution, and the crossover depends on viscosity and the associated high-frequency content~\cite{buitrago2025memory}. Whether an analogous crossover exists for global weather forecasting, and at what grid resolution it might occur, remains an empirical question.

Here, we take a first step toward addressing this question by testing whether temporal memory can improve GraphCast forecasts at $1^\circ$ resolution. We augment a frozen GraphCast model~\cite{lam2023graphcast} with a trainable Mamba correction branch that combines spatial message passing with Mamba layers~\cite{gu2023mamba}. This branch carries a temporal memory state across autoregressive forecast steps and adds a correction to each GraphCast prediction.

Learning a residual correction focuses training on what the pretrained predictor misses. Freezing the predictor leaves only the correction branch's parameters to optimize. The Mori--Zwanzig decomposition offers a structural motivation for this \emph{additive} ansatz: a local-in-time contribution and a history-dependent contribution enter the resolved evolution as separate terms~\cite{YenTingLinMZ2021,de2026data}. The outputs are added linearly, while the learned memory map itself can be nonlinear.

Output corrections to frozen time-series forecasters have been studied in $\delta$-Adapter~\cite{liang2026forecast}, and neural corrections have been applied during the integration of a numerical weather model~\cite{chen2022correcting}. Memory within a forecaster is also established, including recurrent global weather models~\cite{hu2023swinvrnn} and the neural-operator approach above. Our focus is on temporal memory in a separate correction branch added to a frozen pretrained forecaster.

We study the general output-level form
\begin{equation}
    (c_t,m_t)=R_\phi(X_t,m_{t-1}),
    \qquad
    \widehat{Y}_t=G_\theta(X_t)+c_t,
    \qquad \theta\ \text{fixed},
    \label{eq:intro_memory_correction}
\end{equation}
where $G_\theta$ is the pretrained GraphCast backbone, $R_\phi$ is the trainable Mamba correction branch, $c_t$ is its forecast correction, and $m_t$ is its temporal memory state. This state persists between forecast calls and comprises all state-space model (SSM) states and causal-convolution caches in the branch. In our implementation, $G_\theta$ is GraphCast and $R_\phi$ is the correction branch described above. The predictor and correction meet at the output, so this formulation does not require access to the predictor's internal layers. A zero-initialized output head makes the GraphCast--Mamba model start from the GraphCast forecast; only the correction branch is then trained.

We test this instantiation at $1^\circ$ resolution on 32 matched forecast initializations from a development set. The selected correction reduces the GraphCast normalized weighted mean-squared error (GraphCast loss) by $16.34\%$ over ten-day rollouts and $23.68\%$ at day 10 (240\,h after initialization). Separate memory-reset controls, measured with an approximate reconstruction of that loss, support a contribution from persistent temporal context beyond the branch's additional trainable capacity. These results establish a case for memory-based correction of a strong pretrained forecaster at the tested resolution. The resolution at which its memory benefit might diminish, and whether the same correction strategy transfers to other forecasters, require separate matched experiments.
